\documentclass[aspectratio=169,10pt]{beamer}
\makeatletter\def\input@path{{../../estilo/}{../estilo/}{estilo/}}\makeatother
\usepackage{devmobile}
\modulo[Flutter]{19}{Código Nativo, Web, Desktop e Embarcados}{Canais de plataforma, Pigeon, FFI, criação de plugins, Flutter dentro de apps nativos, Flutter web, desktop e Raspberry Pi}

\begin{document}
\capa

\begin{frame}{Objetivos do módulo}
  \begin{conceito}[Ao final deste módulo você será capaz de]
    \begin{itemize}
      \item Comparar as abordagens multiplataforma e justificar uma escolha.
      \item Chamar código Kotlin e Swift a partir do Dart com canais de plataforma e Pigeon.
      \item Usar bibliotecas C com FFI.
      \item Criar um plugin próprio e embutir telas Flutter em apps nativos.
      \item Publicar o app na web e no desktop, adaptando o layout a cada tela.
      \item Executar um painel Flutter num Raspberry Pi.
    \end{itemize}
  \end{conceito}
\end{frame}

\sumario

% =============================================================================
%  Fundamentos compartilhados --- Módulo 19: desenvolvimento multiplataforma
%  Incluído pelas trilhas Kotlin e Flutter com % =============================================================================
%  Fundamentos compartilhados --- Módulo 19: desenvolvimento multiplataforma
%  Incluído pelas trilhas Kotlin e Flutter com % =============================================================================
%  Fundamentos compartilhados --- Módulo 19: desenvolvimento multiplataforma
%  Incluído pelas trilhas Kotlin e Flutter com \input{../fundamentos/f19-multiplataforma}
% =============================================================================
\section{Fundamentos do desenvolvimento multiplataforma}

\begin{frame}{O problema: muitas plataformas, uma equipe}
  \begin{columns}[c]
    \begin{column}{0.5\textwidth}
      \begin{conceito}[Por que existe o multiplataforma]
        Android e iOS têm linguagens, ferramentas e componentes diferentes. Dois apps nativos significam duas bases de código, duas equipes e o dobro de testes para cada funcionalidade.
      \end{conceito}
      \begin{definicao}[Desenvolvimento multiplataforma]
        Escrever uma vez parte (ou todo) o código e executá-lo em mais de um sistema, aceitando algum compromisso em desempenho, aparência ou acesso ao hardware.
      \end{definicao}
    \end{column}
    \begin{column}{0.48\textwidth}
      \begin{historia}
        Java prometia ``escreva uma vez, rode em qualquer lugar'' nos anos 1990. No mobile, vieram os apps híbridos com WebView (PhoneGap/Cordova, 2009), o React Native (2015), o Flutter (2017--2018) e o Kotlin Multiplatform (estável em 2023).
      \end{historia}
      \begin{atencao}
        Nenhuma abordagem elimina o conhecimento de cada plataforma: permissões, publicação, ciclo de vida e integrações nativas continuam diferentes.
      \end{atencao}
    \end{column}
  \end{columns}
\end{frame}

\begin{frame}{O espectro das abordagens}
  \centering
  \begin{tikzpicture}[every node/.style={font=\scriptsize,align=center,minimum height=12mm,text width=2.05cm}]
    \node[caixa] (w) at (0,0) {\textbf{Web / PWA}\\ \tiny HTML, CSS, JS no navegador};
    \node[caixa] (h) at (2.35,0) {\textbf{Híbrido}\\ \tiny site dentro de uma WebView (Capacitor)};
    \node[caixapetroleo] (r) at (4.7,0) {\textbf{Componentes nativos}\\ \tiny React Native};
    \node[caixaambar] (f) at (7.05,0) {\textbf{Motor próprio}\\ \tiny Flutter};
    \node[caixapetroleo] (k) at (9.4,0) {\textbf{Lógica compartilhada}\\ \tiny Kotlin Multiplatform};
    \node[caixaescura] (n) at (11.75,0) {\textbf{Nativo}\\ \tiny Kotlin, Swift};
    \draw[setaambar,<->] (-1.0,-1.0) -- node[rotulo,below,text width=8cm]{\tiny mais código compartilhado \hfill mais fidelidade e acesso nativo} (12.75,-1.0);
  \end{tikzpicture}
  \vspace{4pt}
  \scriptsize
  \begin{tabularx}{\textwidth}{@{}l L L L@{}}
    \toprule
    & \cabecalho{O que se compartilha} & \cabecalho{Quem desenha a tela} & \cabecalho{Compromisso principal} \\
    \midrule
    Web/PWA e híbrido & tudo & o navegador (HTML/CSS) & desempenho e acesso limitado ao hardware \\
    React Native & UI e lógica em JS/TS & componentes nativos de cada sistema & ponte entre JS e o nativo \\
    Flutter & UI e lógica em Dart & o motor do Flutter, pixel a pixel & tamanho do app; aparência própria \\
    KMP & lógica (e opcionalmente a UI) & nativa ou Compose Multiplatform & UI pode exigir código por plataforma \\
    \bottomrule
  \end{tabularx}
\end{frame}

\begin{frame}{Como o código executa em cada abordagem}
  \begin{columns}[T]
    \begin{column}{0.52\textwidth}
      \scriptsize
      \begin{tabularx}{\linewidth}{@{}l L@{}}
        \toprule
        \cabecalho{Tecnologia} & \cabecalho{Execução no aparelho} \\
        \midrule
        PWA / híbrido & JavaScript interpretado e compilado sob demanda (JIT) pelo motor do navegador \\
        React Native & JavaScript pré-compilado em \emph{bytecode} e executado pelo motor Hermes \\
        Flutter & Dart compilado antes (AOT) para código de máquina ARM \\
        KMP (Android) & Kotlin $\rightarrow$ \emph{bytecode} $\rightarrow$ DEX, executado pelo ART \\
        KMP (iOS) & Kotlin/Native compilado via LLVM para binário nativo \\
        Nativo & Kotlin no ART; Swift compilado via LLVM \\
        \bottomrule
      \end{tabularx}
    \end{column}
    \begin{column}{0.46\textwidth}
      \begin{definicao}[JIT e AOT]
        \textbf{JIT} (\emph{just-in-time}): compila durante a execução; flexível, permite recarregar código.
        \textbf{AOT} (\emph{ahead-of-time}): compila antes da distribuição; inicia mais rápido e tem desempenho previsível.
      \end{definicao}
      \begin{pordentro}[Por que o iOS importa aqui]
        A App Store não permite que apps gerem e executem código de máquina novo em tempo de execução (exceto o motor do navegador). Por isso o Flutter usa AOT no iOS e o React Native usa um interpretador de \emph{bytecode}.
      \end{pordentro}
    \end{column}
  \end{columns}
\end{frame}

\begin{frame}{Quem desenha os pixels?}
  \begin{columns}[c]
    \begin{column}{0.56\textwidth}
      \centering
      \begin{tikzpicture}[node distance=2mm,every node/.style={font=\scriptsize,minimum width=2.6cm}]
        \node[caixa] (a1) {Seu código};
        \node[caixapetroleo,below=of a1] (a2) {Componentes nativos\\ \tiny \texttt{UIButton}, \texttt{View}};
        \node[caixaescura,below=of a2] (a3) {Sistema desenha};
        \node[above=1mm of a1,font=\scriptsize\bfseries,text=azulpetroleo] {Nativo / React Native};
        \node[caixa,right=4mm of a1] (b1) {Seu código};
        \node[caixaambar,below=of b1] (b2) {Motor próprio\\ \tiny Impeller / Skia};
        \node[caixaescura,below=of b2] (b3) {GPU (Metal, Vulkan)};
        \node[above=1mm of b1,font=\scriptsize\bfseries,text=azulpetroleo] {Flutter, Compose MP};
        \draw[seta] (a1) -- (a2); \draw[seta] (a2) -- (a3);
        \draw[seta] (b1) -- (b2); \draw[seta] (b2) -- (b3);
      \end{tikzpicture}
    \end{column}
    \begin{column}{0.42\textwidth}
      \begin{conceito}[Componentes nativos]
        A tela usa os controles do sistema: aparência e comportamento idênticos aos de outros apps e atualizados junto com o sistema.
      \end{conceito}
      \begin{conceito}[Motor próprio]
        O framework desenha cada pixel numa tela em branco: aparência idêntica em todas as plataformas e controle total de animações, mas precisa reimplementar o comportamento de cada controle.
      \end{conceito}
      \begin{analogia}
        Montar a casa com móveis de loja (nativos) ou fabricar os próprios móveis (motor próprio).
      \end{analogia}
    \end{column}
  \end{columns}
\end{frame}

\begin{frame}{Conversando com o nativo: pontes, canais e FFI}
  \begin{columns}[T]
    \begin{column}{0.5\textwidth}
      \scriptsize
      \begin{tabularx}{\linewidth}{@{}p{2.2cm} L@{}}
        \toprule
        \cabecalho{Mecanismo} & \cabecalho{Como funciona} \\
        \midrule
        \textbf{Troca de mensagens} & chamadas serializadas e assíncronas entre os dois lados (canais de plataforma do Flutter, ponte antiga do React Native) \\
        \textbf{Interface direta} & chamada de funções nativas sem serializar, inclusive síncrona (FFI do Dart, JSI do React Native) \\
        \textbf{Geração de código} & a interface é declarada uma vez e o código dos dois lados é gerado (Pigeon, Turbo Modules, exportação para Swift do KMP) \\
        \bottomrule
      \end{tabularx}
    \end{column}
    \begin{column}{0.48\textwidth}
      \begin{pordentro}[O custo de atravessar]
        Cada chamada por mensagem converte os dados para bytes, troca de \emph{thread} e converte de volta. Para chamadas raras (abrir a câmera) não importa; para milhares por segundo (áudio, imagens), use interfaces diretas.
      \end{pordentro}
      \begin{definicao}[FFI --- \emph{Foreign Function Interface}]
        Mecanismo para chamar funções escritas em outra linguagem, em geral bibliotecas C, compartilhando memória.
      \end{definicao}
    \end{column}
  \end{columns}
\end{frame}

\begin{frame}{Progressive Web Apps}
  \begin{columns}[T]
    \begin{column}{0.5\textwidth}
      \begin{definicao}[PWA]
        Site que se comporta como app: pode ser instalado na tela inicial, funciona offline e recebe notificações. Base: um \textbf{manifesto} (nome, ícones, cores) e um \textbf{\emph{service worker}} (script que intercepta requisições e gerencia o cache).
      \end{definicao}
      \begin{conceito}[Vantagens]
        Uma única base web, sem loja e sem revisão, atualização instantânea, link compartilhável.
      \end{conceito}
    \end{column}
    \begin{column}{0.48\textwidth}
      \begin{atencao}[Limites]
        \begin{itemize}
          \item Acesso a hardware depende do navegador: Web Bluetooth e Web NFC não existem no Safari.
          \item No iOS, notificações exigem que o usuário instale o site na tela inicial.
          \item Menos presença nas lojas e menos integração com o sistema.
        \end{itemize}
      \end{atencao}
      \begin{dica}
        Para um painel de consulta do LabMonitor no computador do laboratório, uma versão web basta; para BLE e NFC, o app nativo ou multiplataforma é necessário.
      \end{dica}
    \end{column}
  \end{columns}
\end{frame}

\begin{frame}{Decidindo com uma matriz ponderada}
  \begin{columns}[T]
    \begin{column}{0.58\textwidth}
      \scriptsize
      \begin{tabular}{@{}l c c c c c@{}}
        \toprule
        \cabecalho{Critério} & \cabecalho{Peso} & \cabecalho{Nativo} & \cabecalho{Flutter} & \cabecalho{RN} & \cabecalho{KMP} \\
        \midrule
        Atender Android e iOS & 3 & 2 & 5 & 5 & 4 \\
        Acesso a BLE, NFC, câmera & 3 & 5 & 4 & 4 & 5 \\
        Equipe já sabe a linguagem & 2 & 4 & 3 & 2 & 4 \\
        Tempo até a 1ª versão & 2 & 2 & 5 & 4 & 3 \\
        Manutenção em 5 anos & 1 & 4 & 4 & 3 & 4 \\
        \midrule
        \textbf{Total ponderado} & & 37 & \textbf{47} & 42 & 45 \\
        \bottomrule
      \end{tabular}\\[2pt]
      {\tiny Notas de 1 a 5 ilustrativas: app novo, equipe que conhece Kotlin. Com um app Kotlin já existente, o critério ``reaproveitar código'' favoreceria o KMP.}
    \end{column}
    \begin{column}{0.4\textwidth}
      \begin{definicao}[Matriz de decisão ponderada]
        Lista os critérios, dá um peso a cada um conforme sua importância, atribui notas às opções e soma nota $\times$ peso.
      \end{definicao}
      \begin{atencao}
        O resultado depende dos pesos: discuta-os com a equipe e o cliente \textbf{antes} de dar as notas, para não ajustar a matriz ao resultado desejado.
      \end{atencao}
    \end{column}
  \end{columns}
\end{frame}

\begin{frame}{Além do celular: telas e formatos}
  \begin{columns}[T]
    \begin{column}{0.5\textwidth}
      \begin{definicao}[Responsivo e adaptativo]
        \textbf{Responsivo}: o mesmo layout se estica ou reorganiza conforme o espaço.
        \textbf{Adaptativo}: layouts diferentes para classes de tela e para cada forma de entrada (toque, mouse, teclado, controle remoto).
      \end{definicao}
      \scriptsize
      \begin{tabularx}{\linewidth}{@{}l l L@{}}
        \toprule
        \cabecalho{Classe} & \cabecalho{Largura} & \cabecalho{Exemplos} \\
        \midrule
        Compacta & $<$ 600\,dp & celular em pé \\
        Média & 600--839\,dp & tablet em pé, dobrável aberto \\
        Expandida & $\geq$ 840\,dp & tablet deitado, desktop \\
        \bottomrule
      \end{tabularx}\\[2pt]
      {\tiny Classes de janela do Material Design.}
    \end{column}
    \begin{column}{0.48\textwidth}
      \begin{conceito}[Outros formatos]
        Relógios, TVs, carros (Android Auto, CarPlay), painéis embarcados e quiosques. Cada um muda a distância de leitura, a forma de entrada e o tempo de atenção do usuário.
      \end{conceito}
      \begin{dica}
        Projete pelo \textbf{espaço disponível}, não pelo tipo de aparelho: uma janela de desktop pode ser estreita como um celular.
      \end{dica}
    \end{column}
  \end{columns}
\end{frame}

% =============================================================================
\section{Fundamentos do desenvolvimento multiplataforma}

\begin{frame}{O problema: muitas plataformas, uma equipe}
  \begin{columns}[c]
    \begin{column}{0.5\textwidth}
      \begin{conceito}[Por que existe o multiplataforma]
        Android e iOS têm linguagens, ferramentas e componentes diferentes. Dois apps nativos significam duas bases de código, duas equipes e o dobro de testes para cada funcionalidade.
      \end{conceito}
      \begin{definicao}[Desenvolvimento multiplataforma]
        Escrever uma vez parte (ou todo) o código e executá-lo em mais de um sistema, aceitando algum compromisso em desempenho, aparência ou acesso ao hardware.
      \end{definicao}
    \end{column}
    \begin{column}{0.48\textwidth}
      \begin{historia}
        Java prometia ``escreva uma vez, rode em qualquer lugar'' nos anos 1990. No mobile, vieram os apps híbridos com WebView (PhoneGap/Cordova, 2009), o React Native (2015), o Flutter (2017--2018) e o Kotlin Multiplatform (estável em 2023).
      \end{historia}
      \begin{atencao}
        Nenhuma abordagem elimina o conhecimento de cada plataforma: permissões, publicação, ciclo de vida e integrações nativas continuam diferentes.
      \end{atencao}
    \end{column}
  \end{columns}
\end{frame}

\begin{frame}{O espectro das abordagens}
  \centering
  \begin{tikzpicture}[every node/.style={font=\scriptsize,align=center,minimum height=12mm,text width=2.05cm}]
    \node[caixa] (w) at (0,0) {\textbf{Web / PWA}\\ \tiny HTML, CSS, JS no navegador};
    \node[caixa] (h) at (2.35,0) {\textbf{Híbrido}\\ \tiny site dentro de uma WebView (Capacitor)};
    \node[caixapetroleo] (r) at (4.7,0) {\textbf{Componentes nativos}\\ \tiny React Native};
    \node[caixaambar] (f) at (7.05,0) {\textbf{Motor próprio}\\ \tiny Flutter};
    \node[caixapetroleo] (k) at (9.4,0) {\textbf{Lógica compartilhada}\\ \tiny Kotlin Multiplatform};
    \node[caixaescura] (n) at (11.75,0) {\textbf{Nativo}\\ \tiny Kotlin, Swift};
    \draw[setaambar,<->] (-1.0,-1.0) -- node[rotulo,below,text width=8cm]{\tiny mais código compartilhado \hfill mais fidelidade e acesso nativo} (12.75,-1.0);
  \end{tikzpicture}
  \vspace{4pt}
  \scriptsize
  \begin{tabularx}{\textwidth}{@{}l L L L@{}}
    \toprule
    & \cabecalho{O que se compartilha} & \cabecalho{Quem desenha a tela} & \cabecalho{Compromisso principal} \\
    \midrule
    Web/PWA e híbrido & tudo & o navegador (HTML/CSS) & desempenho e acesso limitado ao hardware \\
    React Native & UI e lógica em JS/TS & componentes nativos de cada sistema & ponte entre JS e o nativo \\
    Flutter & UI e lógica em Dart & o motor do Flutter, pixel a pixel & tamanho do app; aparência própria \\
    KMP & lógica (e opcionalmente a UI) & nativa ou Compose Multiplatform & UI pode exigir código por plataforma \\
    \bottomrule
  \end{tabularx}
\end{frame}

\begin{frame}{Como o código executa em cada abordagem}
  \begin{columns}[T]
    \begin{column}{0.52\textwidth}
      \scriptsize
      \begin{tabularx}{\linewidth}{@{}l L@{}}
        \toprule
        \cabecalho{Tecnologia} & \cabecalho{Execução no aparelho} \\
        \midrule
        PWA / híbrido & JavaScript interpretado e compilado sob demanda (JIT) pelo motor do navegador \\
        React Native & JavaScript pré-compilado em \emph{bytecode} e executado pelo motor Hermes \\
        Flutter & Dart compilado antes (AOT) para código de máquina ARM \\
        KMP (Android) & Kotlin $\rightarrow$ \emph{bytecode} $\rightarrow$ DEX, executado pelo ART \\
        KMP (iOS) & Kotlin/Native compilado via LLVM para binário nativo \\
        Nativo & Kotlin no ART; Swift compilado via LLVM \\
        \bottomrule
      \end{tabularx}
    \end{column}
    \begin{column}{0.46\textwidth}
      \begin{definicao}[JIT e AOT]
        \textbf{JIT} (\emph{just-in-time}): compila durante a execução; flexível, permite recarregar código.
        \textbf{AOT} (\emph{ahead-of-time}): compila antes da distribuição; inicia mais rápido e tem desempenho previsível.
      \end{definicao}
      \begin{pordentro}[Por que o iOS importa aqui]
        A App Store não permite que apps gerem e executem código de máquina novo em tempo de execução (exceto o motor do navegador). Por isso o Flutter usa AOT no iOS e o React Native usa um interpretador de \emph{bytecode}.
      \end{pordentro}
    \end{column}
  \end{columns}
\end{frame}

\begin{frame}{Quem desenha os pixels?}
  \begin{columns}[c]
    \begin{column}{0.56\textwidth}
      \centering
      \begin{tikzpicture}[node distance=2mm,every node/.style={font=\scriptsize,minimum width=2.6cm}]
        \node[caixa] (a1) {Seu código};
        \node[caixapetroleo,below=of a1] (a2) {Componentes nativos\\ \tiny \texttt{UIButton}, \texttt{View}};
        \node[caixaescura,below=of a2] (a3) {Sistema desenha};
        \node[above=1mm of a1,font=\scriptsize\bfseries,text=azulpetroleo] {Nativo / React Native};
        \node[caixa,right=4mm of a1] (b1) {Seu código};
        \node[caixaambar,below=of b1] (b2) {Motor próprio\\ \tiny Impeller / Skia};
        \node[caixaescura,below=of b2] (b3) {GPU (Metal, Vulkan)};
        \node[above=1mm of b1,font=\scriptsize\bfseries,text=azulpetroleo] {Flutter, Compose MP};
        \draw[seta] (a1) -- (a2); \draw[seta] (a2) -- (a3);
        \draw[seta] (b1) -- (b2); \draw[seta] (b2) -- (b3);
      \end{tikzpicture}
    \end{column}
    \begin{column}{0.42\textwidth}
      \begin{conceito}[Componentes nativos]
        A tela usa os controles do sistema: aparência e comportamento idênticos aos de outros apps e atualizados junto com o sistema.
      \end{conceito}
      \begin{conceito}[Motor próprio]
        O framework desenha cada pixel numa tela em branco: aparência idêntica em todas as plataformas e controle total de animações, mas precisa reimplementar o comportamento de cada controle.
      \end{conceito}
      \begin{analogia}
        Montar a casa com móveis de loja (nativos) ou fabricar os próprios móveis (motor próprio).
      \end{analogia}
    \end{column}
  \end{columns}
\end{frame}

\begin{frame}{Conversando com o nativo: pontes, canais e FFI}
  \begin{columns}[T]
    \begin{column}{0.5\textwidth}
      \scriptsize
      \begin{tabularx}{\linewidth}{@{}p{2.2cm} L@{}}
        \toprule
        \cabecalho{Mecanismo} & \cabecalho{Como funciona} \\
        \midrule
        \textbf{Troca de mensagens} & chamadas serializadas e assíncronas entre os dois lados (canais de plataforma do Flutter, ponte antiga do React Native) \\
        \textbf{Interface direta} & chamada de funções nativas sem serializar, inclusive síncrona (FFI do Dart, JSI do React Native) \\
        \textbf{Geração de código} & a interface é declarada uma vez e o código dos dois lados é gerado (Pigeon, Turbo Modules, exportação para Swift do KMP) \\
        \bottomrule
      \end{tabularx}
    \end{column}
    \begin{column}{0.48\textwidth}
      \begin{pordentro}[O custo de atravessar]
        Cada chamada por mensagem converte os dados para bytes, troca de \emph{thread} e converte de volta. Para chamadas raras (abrir a câmera) não importa; para milhares por segundo (áudio, imagens), use interfaces diretas.
      \end{pordentro}
      \begin{definicao}[FFI --- \emph{Foreign Function Interface}]
        Mecanismo para chamar funções escritas em outra linguagem, em geral bibliotecas C, compartilhando memória.
      \end{definicao}
    \end{column}
  \end{columns}
\end{frame}

\begin{frame}{Progressive Web Apps}
  \begin{columns}[T]
    \begin{column}{0.5\textwidth}
      \begin{definicao}[PWA]
        Site que se comporta como app: pode ser instalado na tela inicial, funciona offline e recebe notificações. Base: um \textbf{manifesto} (nome, ícones, cores) e um \textbf{\emph{service worker}} (script que intercepta requisições e gerencia o cache).
      \end{definicao}
      \begin{conceito}[Vantagens]
        Uma única base web, sem loja e sem revisão, atualização instantânea, link compartilhável.
      \end{conceito}
    \end{column}
    \begin{column}{0.48\textwidth}
      \begin{atencao}[Limites]
        \begin{itemize}
          \item Acesso a hardware depende do navegador: Web Bluetooth e Web NFC não existem no Safari.
          \item No iOS, notificações exigem que o usuário instale o site na tela inicial.
          \item Menos presença nas lojas e menos integração com o sistema.
        \end{itemize}
      \end{atencao}
      \begin{dica}
        Para um painel de consulta do LabMonitor no computador do laboratório, uma versão web basta; para BLE e NFC, o app nativo ou multiplataforma é necessário.
      \end{dica}
    \end{column}
  \end{columns}
\end{frame}

\begin{frame}{Decidindo com uma matriz ponderada}
  \begin{columns}[T]
    \begin{column}{0.58\textwidth}
      \scriptsize
      \begin{tabular}{@{}l c c c c c@{}}
        \toprule
        \cabecalho{Critério} & \cabecalho{Peso} & \cabecalho{Nativo} & \cabecalho{Flutter} & \cabecalho{RN} & \cabecalho{KMP} \\
        \midrule
        Atender Android e iOS & 3 & 2 & 5 & 5 & 4 \\
        Acesso a BLE, NFC, câmera & 3 & 5 & 4 & 4 & 5 \\
        Equipe já sabe a linguagem & 2 & 4 & 3 & 2 & 4 \\
        Tempo até a 1ª versão & 2 & 2 & 5 & 4 & 3 \\
        Manutenção em 5 anos & 1 & 4 & 4 & 3 & 4 \\
        \midrule
        \textbf{Total ponderado} & & 37 & \textbf{47} & 42 & 45 \\
        \bottomrule
      \end{tabular}\\[2pt]
      {\tiny Notas de 1 a 5 ilustrativas: app novo, equipe que conhece Kotlin. Com um app Kotlin já existente, o critério ``reaproveitar código'' favoreceria o KMP.}
    \end{column}
    \begin{column}{0.4\textwidth}
      \begin{definicao}[Matriz de decisão ponderada]
        Lista os critérios, dá um peso a cada um conforme sua importância, atribui notas às opções e soma nota $\times$ peso.
      \end{definicao}
      \begin{atencao}
        O resultado depende dos pesos: discuta-os com a equipe e o cliente \textbf{antes} de dar as notas, para não ajustar a matriz ao resultado desejado.
      \end{atencao}
    \end{column}
  \end{columns}
\end{frame}

\begin{frame}{Além do celular: telas e formatos}
  \begin{columns}[T]
    \begin{column}{0.5\textwidth}
      \begin{definicao}[Responsivo e adaptativo]
        \textbf{Responsivo}: o mesmo layout se estica ou reorganiza conforme o espaço.
        \textbf{Adaptativo}: layouts diferentes para classes de tela e para cada forma de entrada (toque, mouse, teclado, controle remoto).
      \end{definicao}
      \scriptsize
      \begin{tabularx}{\linewidth}{@{}l l L@{}}
        \toprule
        \cabecalho{Classe} & \cabecalho{Largura} & \cabecalho{Exemplos} \\
        \midrule
        Compacta & $<$ 600\,dp & celular em pé \\
        Média & 600--839\,dp & tablet em pé, dobrável aberto \\
        Expandida & $\geq$ 840\,dp & tablet deitado, desktop \\
        \bottomrule
      \end{tabularx}\\[2pt]
      {\tiny Classes de janela do Material Design.}
    \end{column}
    \begin{column}{0.48\textwidth}
      \begin{conceito}[Outros formatos]
        Relógios, TVs, carros (Android Auto, CarPlay), painéis embarcados e quiosques. Cada um muda a distância de leitura, a forma de entrada e o tempo de atenção do usuário.
      \end{conceito}
      \begin{dica}
        Projete pelo \textbf{espaço disponível}, não pelo tipo de aparelho: uma janela de desktop pode ser estreita como um celular.
      \end{dica}
    \end{column}
  \end{columns}
\end{frame}

% =============================================================================
\section{Fundamentos do desenvolvimento multiplataforma}

\begin{frame}{O problema: muitas plataformas, uma equipe}
  \begin{columns}[c]
    \begin{column}{0.5\textwidth}
      \begin{conceito}[Por que existe o multiplataforma]
        Android e iOS têm linguagens, ferramentas e componentes diferentes. Dois apps nativos significam duas bases de código, duas equipes e o dobro de testes para cada funcionalidade.
      \end{conceito}
      \begin{definicao}[Desenvolvimento multiplataforma]
        Escrever uma vez parte (ou todo) o código e executá-lo em mais de um sistema, aceitando algum compromisso em desempenho, aparência ou acesso ao hardware.
      \end{definicao}
    \end{column}
    \begin{column}{0.48\textwidth}
      \begin{historia}
        Java prometia ``escreva uma vez, rode em qualquer lugar'' nos anos 1990. No mobile, vieram os apps híbridos com WebView (PhoneGap/Cordova, 2009), o React Native (2015), o Flutter (2017--2018) e o Kotlin Multiplatform (estável em 2023).
      \end{historia}
      \begin{atencao}
        Nenhuma abordagem elimina o conhecimento de cada plataforma: permissões, publicação, ciclo de vida e integrações nativas continuam diferentes.
      \end{atencao}
    \end{column}
  \end{columns}
\end{frame}

\begin{frame}{O espectro das abordagens}
  \centering
  \begin{tikzpicture}[every node/.style={font=\scriptsize,align=center,minimum height=12mm,text width=2.05cm}]
    \node[caixa] (w) at (0,0) {\textbf{Web / PWA}\\ \tiny HTML, CSS, JS no navegador};
    \node[caixa] (h) at (2.35,0) {\textbf{Híbrido}\\ \tiny site dentro de uma WebView (Capacitor)};
    \node[caixapetroleo] (r) at (4.7,0) {\textbf{Componentes nativos}\\ \tiny React Native};
    \node[caixaambar] (f) at (7.05,0) {\textbf{Motor próprio}\\ \tiny Flutter};
    \node[caixapetroleo] (k) at (9.4,0) {\textbf{Lógica compartilhada}\\ \tiny Kotlin Multiplatform};
    \node[caixaescura] (n) at (11.75,0) {\textbf{Nativo}\\ \tiny Kotlin, Swift};
    \draw[setaambar,<->] (-1.0,-1.0) -- node[rotulo,below,text width=8cm]{\tiny mais código compartilhado \hfill mais fidelidade e acesso nativo} (12.75,-1.0);
  \end{tikzpicture}
  \vspace{4pt}
  \scriptsize
  \begin{tabularx}{\textwidth}{@{}l L L L@{}}
    \toprule
    & \cabecalho{O que se compartilha} & \cabecalho{Quem desenha a tela} & \cabecalho{Compromisso principal} \\
    \midrule
    Web/PWA e híbrido & tudo & o navegador (HTML/CSS) & desempenho e acesso limitado ao hardware \\
    React Native & UI e lógica em JS/TS & componentes nativos de cada sistema & ponte entre JS e o nativo \\
    Flutter & UI e lógica em Dart & o motor do Flutter, pixel a pixel & tamanho do app; aparência própria \\
    KMP & lógica (e opcionalmente a UI) & nativa ou Compose Multiplatform & UI pode exigir código por plataforma \\
    \bottomrule
  \end{tabularx}
\end{frame}

\begin{frame}{Como o código executa em cada abordagem}
  \begin{columns}[T]
    \begin{column}{0.52\textwidth}
      \scriptsize
      \begin{tabularx}{\linewidth}{@{}l L@{}}
        \toprule
        \cabecalho{Tecnologia} & \cabecalho{Execução no aparelho} \\
        \midrule
        PWA / híbrido & JavaScript interpretado e compilado sob demanda (JIT) pelo motor do navegador \\
        React Native & JavaScript pré-compilado em \emph{bytecode} e executado pelo motor Hermes \\
        Flutter & Dart compilado antes (AOT) para código de máquina ARM \\
        KMP (Android) & Kotlin $\rightarrow$ \emph{bytecode} $\rightarrow$ DEX, executado pelo ART \\
        KMP (iOS) & Kotlin/Native compilado via LLVM para binário nativo \\
        Nativo & Kotlin no ART; Swift compilado via LLVM \\
        \bottomrule
      \end{tabularx}
    \end{column}
    \begin{column}{0.46\textwidth}
      \begin{definicao}[JIT e AOT]
        \textbf{JIT} (\emph{just-in-time}): compila durante a execução; flexível, permite recarregar código.
        \textbf{AOT} (\emph{ahead-of-time}): compila antes da distribuição; inicia mais rápido e tem desempenho previsível.
      \end{definicao}
      \begin{pordentro}[Por que o iOS importa aqui]
        A App Store não permite que apps gerem e executem código de máquina novo em tempo de execução (exceto o motor do navegador). Por isso o Flutter usa AOT no iOS e o React Native usa um interpretador de \emph{bytecode}.
      \end{pordentro}
    \end{column}
  \end{columns}
\end{frame}

\begin{frame}{Quem desenha os pixels?}
  \begin{columns}[c]
    \begin{column}{0.56\textwidth}
      \centering
      \begin{tikzpicture}[node distance=2mm,every node/.style={font=\scriptsize,minimum width=2.6cm}]
        \node[caixa] (a1) {Seu código};
        \node[caixapetroleo,below=of a1] (a2) {Componentes nativos\\ \tiny \texttt{UIButton}, \texttt{View}};
        \node[caixaescura,below=of a2] (a3) {Sistema desenha};
        \node[above=1mm of a1,font=\scriptsize\bfseries,text=azulpetroleo] {Nativo / React Native};
        \node[caixa,right=4mm of a1] (b1) {Seu código};
        \node[caixaambar,below=of b1] (b2) {Motor próprio\\ \tiny Impeller / Skia};
        \node[caixaescura,below=of b2] (b3) {GPU (Metal, Vulkan)};
        \node[above=1mm of b1,font=\scriptsize\bfseries,text=azulpetroleo] {Flutter, Compose MP};
        \draw[seta] (a1) -- (a2); \draw[seta] (a2) -- (a3);
        \draw[seta] (b1) -- (b2); \draw[seta] (b2) -- (b3);
      \end{tikzpicture}
    \end{column}
    \begin{column}{0.42\textwidth}
      \begin{conceito}[Componentes nativos]
        A tela usa os controles do sistema: aparência e comportamento idênticos aos de outros apps e atualizados junto com o sistema.
      \end{conceito}
      \begin{conceito}[Motor próprio]
        O framework desenha cada pixel numa tela em branco: aparência idêntica em todas as plataformas e controle total de animações, mas precisa reimplementar o comportamento de cada controle.
      \end{conceito}
      \begin{analogia}
        Montar a casa com móveis de loja (nativos) ou fabricar os próprios móveis (motor próprio).
      \end{analogia}
    \end{column}
  \end{columns}
\end{frame}

\begin{frame}{Conversando com o nativo: pontes, canais e FFI}
  \begin{columns}[T]
    \begin{column}{0.5\textwidth}
      \scriptsize
      \begin{tabularx}{\linewidth}{@{}p{2.2cm} L@{}}
        \toprule
        \cabecalho{Mecanismo} & \cabecalho{Como funciona} \\
        \midrule
        \textbf{Troca de mensagens} & chamadas serializadas e assíncronas entre os dois lados (canais de plataforma do Flutter, ponte antiga do React Native) \\
        \textbf{Interface direta} & chamada de funções nativas sem serializar, inclusive síncrona (FFI do Dart, JSI do React Native) \\
        \textbf{Geração de código} & a interface é declarada uma vez e o código dos dois lados é gerado (Pigeon, Turbo Modules, exportação para Swift do KMP) \\
        \bottomrule
      \end{tabularx}
    \end{column}
    \begin{column}{0.48\textwidth}
      \begin{pordentro}[O custo de atravessar]
        Cada chamada por mensagem converte os dados para bytes, troca de \emph{thread} e converte de volta. Para chamadas raras (abrir a câmera) não importa; para milhares por segundo (áudio, imagens), use interfaces diretas.
      \end{pordentro}
      \begin{definicao}[FFI --- \emph{Foreign Function Interface}]
        Mecanismo para chamar funções escritas em outra linguagem, em geral bibliotecas C, compartilhando memória.
      \end{definicao}
    \end{column}
  \end{columns}
\end{frame}

\begin{frame}{Progressive Web Apps}
  \begin{columns}[T]
    \begin{column}{0.5\textwidth}
      \begin{definicao}[PWA]
        Site que se comporta como app: pode ser instalado na tela inicial, funciona offline e recebe notificações. Base: um \textbf{manifesto} (nome, ícones, cores) e um \textbf{\emph{service worker}} (script que intercepta requisições e gerencia o cache).
      \end{definicao}
      \begin{conceito}[Vantagens]
        Uma única base web, sem loja e sem revisão, atualização instantânea, link compartilhável.
      \end{conceito}
    \end{column}
    \begin{column}{0.48\textwidth}
      \begin{atencao}[Limites]
        \begin{itemize}
          \item Acesso a hardware depende do navegador: Web Bluetooth e Web NFC não existem no Safari.
          \item No iOS, notificações exigem que o usuário instale o site na tela inicial.
          \item Menos presença nas lojas e menos integração com o sistema.
        \end{itemize}
      \end{atencao}
      \begin{dica}
        Para um painel de consulta do LabMonitor no computador do laboratório, uma versão web basta; para BLE e NFC, o app nativo ou multiplataforma é necessário.
      \end{dica}
    \end{column}
  \end{columns}
\end{frame}

\begin{frame}{Decidindo com uma matriz ponderada}
  \begin{columns}[T]
    \begin{column}{0.58\textwidth}
      \scriptsize
      \begin{tabular}{@{}l c c c c c@{}}
        \toprule
        \cabecalho{Critério} & \cabecalho{Peso} & \cabecalho{Nativo} & \cabecalho{Flutter} & \cabecalho{RN} & \cabecalho{KMP} \\
        \midrule
        Atender Android e iOS & 3 & 2 & 5 & 5 & 4 \\
        Acesso a BLE, NFC, câmera & 3 & 5 & 4 & 4 & 5 \\
        Equipe já sabe a linguagem & 2 & 4 & 3 & 2 & 4 \\
        Tempo até a 1ª versão & 2 & 2 & 5 & 4 & 3 \\
        Manutenção em 5 anos & 1 & 4 & 4 & 3 & 4 \\
        \midrule
        \textbf{Total ponderado} & & 37 & \textbf{47} & 42 & 45 \\
        \bottomrule
      \end{tabular}\\[2pt]
      {\tiny Notas de 1 a 5 ilustrativas: app novo, equipe que conhece Kotlin. Com um app Kotlin já existente, o critério ``reaproveitar código'' favoreceria o KMP.}
    \end{column}
    \begin{column}{0.4\textwidth}
      \begin{definicao}[Matriz de decisão ponderada]
        Lista os critérios, dá um peso a cada um conforme sua importância, atribui notas às opções e soma nota $\times$ peso.
      \end{definicao}
      \begin{atencao}
        O resultado depende dos pesos: discuta-os com a equipe e o cliente \textbf{antes} de dar as notas, para não ajustar a matriz ao resultado desejado.
      \end{atencao}
    \end{column}
  \end{columns}
\end{frame}

\begin{frame}{Além do celular: telas e formatos}
  \begin{columns}[T]
    \begin{column}{0.5\textwidth}
      \begin{definicao}[Responsivo e adaptativo]
        \textbf{Responsivo}: o mesmo layout se estica ou reorganiza conforme o espaço.
        \textbf{Adaptativo}: layouts diferentes para classes de tela e para cada forma de entrada (toque, mouse, teclado, controle remoto).
      \end{definicao}
      \scriptsize
      \begin{tabularx}{\linewidth}{@{}l l L@{}}
        \toprule
        \cabecalho{Classe} & \cabecalho{Largura} & \cabecalho{Exemplos} \\
        \midrule
        Compacta & $<$ 600\,dp & celular em pé \\
        Média & 600--839\,dp & tablet em pé, dobrável aberto \\
        Expandida & $\geq$ 840\,dp & tablet deitado, desktop \\
        \bottomrule
      \end{tabularx}\\[2pt]
      {\tiny Classes de janela do Material Design.}
    \end{column}
    \begin{column}{0.48\textwidth}
      \begin{conceito}[Outros formatos]
        Relógios, TVs, carros (Android Auto, CarPlay), painéis embarcados e quiosques. Cada um muda a distância de leitura, a forma de entrada e o tempo de atenção do usuário.
      \end{conceito}
      \begin{dica}
        Projete pelo \textbf{espaço disponível}, não pelo tipo de aparelho: uma janela de desktop pode ser estreita como um celular.
      \end{dica}
    \end{column}
  \end{columns}
\end{frame}


% =============================================================================
\section{A arquitetura do Flutter}

\begin{frame}{Camadas: framework, motor e \emph{embedder}}
  \begin{columns}[c]
    \begin{column}{0.5\textwidth}
      \centering
      \begin{tikzpicture}[node distance=1.5mm,every node/.style={font=\scriptsize,minimum width=5.2cm}]
        \node[caixa] (app) {Seu app (widgets)};
        \node[caixapetroleo,below=of app] (fw) {\textbf{Framework} (Dart)\\ \tiny Material, Cupertino, widgets, renderização, animação};
        \node[caixaambar,below=of fw] (en) {\textbf{Engine} (C++)\\ \tiny Impeller, Dart runtime, texto, canais};
        \node[caixaescura,below=of en] (em) {\textbf{Embedder} (por plataforma)\\ \tiny Android, iOS, web, Windows, macOS, Linux};
        \node[caixa,below=of em,fill=cinzaclaro] (so) {Sistema operacional e GPU};
      \end{tikzpicture}
    \end{column}
    \begin{column}{0.48\textwidth}
      \begin{conceito}[Quem faz o quê]
        \textbf{Framework}: tudo que você usa em Dart.\\
        \textbf{Engine}: desenha, roda o Dart, faz o \emph{layout} de texto e transporta as mensagens dos canais.\\
        \textbf{Embedder}: cria a janela ou superfície, entrega toques e teclado e integra com o ciclo de vida de cada sistema.
      \end{conceito}
      \begin{pordentro}[Portabilidade]
        Para levar o Flutter a uma plataforma nova, basta escrever um \emph{embedder}. É assim que ele roda em painéis de carros e no Raspberry Pi.
      \end{pordentro}
    \end{column}
  \end{columns}
\end{frame}

% =============================================================================
\section{Canais de plataforma}

\begin{frame}[fragile]{\cod{MethodChannel}: o lado Dart}
  \begin{columns}[T]
    \begin{column}{0.54\textwidth}
\begin{codigo}[title=lib/data/services/bateria\_nativa.dart]{Dart}
import 'package:flutter/services.dart';

class BateriaNativa {
  static const _canal =
      MethodChannel('br.edu.ifma.labmonitor/bateria');

  Future<int> nivel() async {
    try {
      final n = await _canal.invokeMethod<int>('nivel');
      return n ?? -1;
    } on PlatformException catch (e) {
      throw BateriaIndisponivel(e.message);
    } on MissingPluginException {
      return -1;            // plataforma sem implementação
    }
  }
}
\end{codigo}
    \end{column}
    \begin{column}{0.44\textwidth}
      \centering
      \begin{tikzpicture}[node distance=6mm,every node/.style={font=\scriptsize}]
        \node[caixaambar,minimum width=1.8cm] (d) {Dart};
        \node[caixaescura,right=of d,minimum width=1.8cm,align=center] (n) {Kotlin\\ ou Swift};
        \draw[seta] ([yshift=2mm]d.east) -- node[rotulo,above]{\tiny \texttt{'nivel'}} ([yshift=2mm]n.west);
        \draw[setaambar] ([yshift=-2mm]n.west) -- node[rotulo,below]{\tiny 87} ([yshift=-2mm]d.east);
      \end{tikzpicture}
      \begin{conceito}[Canal]
        Nome único + \emph{codec}. A mensagem (nome do método e argumentos) é convertida em bytes, atravessa para o nativo e a resposta volta como \cod{Future}.
      \end{conceito}
      \begin{dica}
        Tipos suportados pelo \emph{codec} padrão: \cod{null}, \cod{bool}, números, \cod{String}, listas, mapas e \cod{Uint8List}.
      \end{dica}
    \end{column}
  \end{columns}
\end{frame}

\begin{frame}[fragile]{\cod{MethodChannel}: os lados Kotlin e Swift}
  \begin{columns}[T]
    \begin{column}{0.52\textwidth}
\begin{codigo}[title=android/.../MainActivity.kt]{Kotlin}
class MainActivity : FlutterActivity() {
  override fun configureFlutterEngine(
      engine: FlutterEngine) {
    super.configureFlutterEngine(engine)
    MethodChannel(engine.dartExecutor.binaryMessenger,
        "br.edu.ifma.labmonitor/bateria")
      .setMethodCallHandler { call, result ->
        if (call.method == "nivel") {
          val bm = getSystemService(BATTERY_SERVICE)
              as BatteryManager
          result.success(bm.getIntProperty(
              BatteryManager.BATTERY_PROPERTY_CAPACITY))
        } else result.notImplemented()
      }
  }
}
\end{codigo}
    \end{column}
    \begin{column}{0.46\textwidth}
\begin{codigo}[title=ios/Runner/CanalBateria.swift]{Swift}
func registrarCanalBateria(
    _ mensageiro: FlutterBinaryMessenger) {
  let canal = FlutterMethodChannel(
    name: "br.edu.ifma.labmonitor/bateria",
    binaryMessenger: mensageiro)
  canal.setMethodCallHandler { call, result in
    guard call.method == "nivel" else {
      return result(FlutterMethodNotImplemented)
    }
    UIDevice.current.isBatteryMonitoringEnabled = true
    result(Int(UIDevice.current.batteryLevel * 100))
  }
}
\end{codigo}
      \begin{dica}
        Registre o canal onde o \cod{AppDelegate} obtém o \emph{messenger} do motor.
      \end{dica}
    \end{column}
  \end{columns}
\end{frame}

\begin{frame}[fragile]{\cod{EventChannel} e Pigeon}
  \begin{columns}[T]
    \begin{column}{0.46\textwidth}
\begin{codigo}[title=Fluxo contínuo do nativo]{Dart}
const _eventos = EventChannel(
    'br.edu.ifma.labmonitor/luz');

Stream<double> luminosidade() =>
    _eventos.receiveBroadcastStream()
        .map((v) => (v as num).toDouble());
\end{codigo}
      \begin{conceito}[\cod{EventChannel}]
        O nativo envia valores quando quiser (sensor, conexão); no Dart, chegam como \cod{Stream}.
      \end{conceito}
    \end{column}
    \begin{column}{0.52\textwidth}
\begin{codigo}[title=pigeons/sensor\_api.dart]{Dart}
import 'package:pigeon/pigeon.dart';
class Leitura {
  Leitura({required this.valor, required this.emMs});
  double valor;
  int emMs;
}

@HostApi()
abstract class SensorNativoApi {
  @async
  Leitura ler(String sensor);
}
\end{codigo}
\begin{terminal}
$ dart run pigeon --input pigeons/sensor_api.dart
\end{terminal}
      \begin{dica}[Pigeon]
        Gera Dart, Kotlin e Swift tipados a partir da interface.
      \end{dica}
    \end{column}
  \end{columns}
\end{frame}

% =============================================================================
\section{FFI: bibliotecas em C}

\begin{frame}[fragile]{Chamando C a partir do Dart}
  \begin{columns}[T]
    \begin{column}{0.36\textwidth}
\begin{codigo}[title=src/filtro.c]{C}
double media_movel(const double *v,
                   int n) {
  double s = 0;
  for (int i = 0; i < n; i++)
    s += v[i];
  return n > 0 ? s / n : 0;
}
\end{codigo}
      \begin{conceito}[Quando usar]
        Reaproveitar bibliotecas C/C++ ou acelerar cálculos chamados com muita frequência.
      \end{conceito}
      \begin{dica}
        \cod{ffigen} gera as assinaturas a partir do \cod{.h}.
      \end{dica}
    \end{column}
    \begin{column}{0.62\textwidth}
\begin{codigo}[title=lib/utils/filtro.dart]{Dart}
import 'dart:ffi';
import 'package:ffi/ffi.dart';
typedef _MediaC = Double Function(Pointer<Double>, Int32);
typedef _MediaDart = double Function(Pointer<Double>, int);

final _lib = DynamicLibrary.open('libfiltro.so');  // Android, Linux
final _media = _lib.lookupFunction<_MediaC, _MediaDart>('media_movel');

double mediaMovel(List<double> valores) {
  final p = calloc<Double>(valores.length);   // memória nativa
  try {
    for (var i = 0; i < valores.length; i++) p[i] = valores[i];
    return _media(p, valores.length);
  } finally {
    calloc.free(p);                           // o GC não libera
  }
}
\end{codigo}
      \begin{atencao}
        Memória nativa não é coletada pelo GC: libere sempre.
      \end{atencao}
    \end{column}
  \end{columns}
\end{frame}

% =============================================================================
\section{Plugins e Flutter em apps nativos}

\begin{frame}[fragile]{Criando um plugin}
  \begin{columns}[T]
    \begin{column}{0.48\textwidth}
\begin{terminal}
$ flutter create --template=plugin \
    --platforms=android,ios \
    --org br.edu.ifma labmonitor_sensor
\end{terminal}
      \centering
      \begin{tikzpicture}[node distance=2mm,every node/.style={font=\scriptsize}]
        \node[caixapetroleo,minimum width=4.6cm] (a) {\cod{labmonitor\_sensor}\\ \tiny API usada pelos apps};
        \node[caixa,minimum width=4.6cm,below=of a] (i) {\cod{labmonitor\_sensor\_platform\_interface}\\ \tiny contrato comum};
        \node[caixaambar,minimum width=1.45cm,below=4mm of i.south west,anchor=north west] (an) {\tiny \_android};
        \node[caixaambar,minimum width=1.45cm,below=4mm of i] (io) {\tiny \_ios};
        \node[caixaambar,minimum width=1.45cm,below=4mm of i.south east,anchor=north east] (w) {\tiny \_web};
        \draw[seta] (a) -- (i);
        \draw[seta] (i.south -| an) -- (an); \draw[seta] (i) -- (io); \draw[seta] (i.south -| w) -- (w);
      \end{tikzpicture}
    \end{column}
    \begin{column}{0.5\textwidth}
      \begin{definicao}[Plugin federado]
        Divide o plugin em pacotes: a API pública, uma interface de plataforma e uma implementação por plataforma. Outras pessoas podem implementar plataformas novas sem mudar o pacote principal.
      \end{definicao}
      \begin{conceito}[Plugin $\times$ pacote]
        \textbf{Pacote} contém só Dart. \textbf{Plugin} tem também código nativo para uma ou mais plataformas.
      \end{conceito}
      \begin{dica}
        Antes de criar um plugin, procure no pub.dev: provavelmente alguém já resolveu o problema.
      \end{dica}
    \end{column}
  \end{columns}
\end{frame}

\begin{frame}[fragile]{\emph{Add-to-app}: Flutter dentro do app Kotlin}
  \begin{columns}[T]
    \begin{column}{0.56\textwidth}
\begin{terminal}
$ flutter create --template=module relatorios_flutter
\end{terminal}
\begin{codigo}[title=App Android nativo (Kotlin)]{Kotlin}
class LabApp : Application() {
  override fun onCreate() {
    super.onCreate()
    // motor pré-aquecido: a tela abre sem atraso
    val motor = FlutterEngine(this).apply {
      dartExecutor.executeDartEntrypoint(
        DartExecutor.DartEntrypoint.createDefault())
    }
    FlutterEngineCache.getInstance().put("relatorios", motor)
  }
}

// ao tocar em "Relatórios"
startActivity(FlutterActivity
    .withCachedEngine("relatorios").build(this))
\end{codigo}
    \end{column}
    \begin{column}{0.42\textwidth}
      \begin{definicao}[\emph{Add-to-app}]
        Incorporar telas Flutter a um app nativo existente, como módulo, em vez de reescrevê-lo inteiro.
      \end{definicao}
      \begin{conceito}[Cenário]
        O LabMonitor em Kotlin ganha a tela de relatórios em Flutter, compartilhada com o iOS (onde se usa um \cod{FlutterViewController}).
      \end{conceito}
      \begin{atencao}
        Cada motor ocupa memória: reaproveite motores em cache ou use \cod{FlutterEngineGroup}.
      \end{atencao}
    \end{column}
  \end{columns}
\end{frame}

% =============================================================================
\section{Flutter na web}

\begin{frame}[fragile]{Compilando para o navegador}
  \begin{columns}[T]
    \begin{column}{0.5\textwidth}
\begin{terminal}
$ flutter run -d chrome
$ flutter build web --wasm --release
# publicar num subcaminho (GitHub Pages)
$ flutter build web --base-href /LabMonitor/
\end{terminal}
      \begin{conceito}[Renderizadores]
        \textbf{CanvasKit}: o Skia compilado para WebAssembly, desenhando num \emph{canvas}.\\
        \textbf{skwasm}: versão mais rápida, usada com \cod{--wasm} quando o navegador suporta WebAssembly com coleta de lixo.
      \end{conceito}
    \end{column}
    \begin{column}{0.48\textwidth}
      \begin{atencao}[Quando não usar]
        Sites de conteúdo que dependem de busca (SEO), como blogs e páginas institucionais: o Flutter web desenha num \emph{canvas} e baixa um motor inteiro. Ele brilha em \textbf{aplicações}: painéis, ferramentas internas, versões web de apps.
      \end{atencao}
      \begin{dica}
        Hospedagem simples: Firebase Hosting, GitHub Pages ou qualquer servidor de arquivos estáticos.
      \end{dica}
    \end{column}
  \end{columns}
\end{frame}

\begin{frame}[fragile]{Código que roda em todo lugar}
  \begin{columns}[T]
    \begin{column}{0.54\textwidth}
\begin{codigo}[title=Diferenças de plataforma]{Dart}
import 'package:flutter/foundation.dart';
// importação condicional: escolhe o arquivo na compilação
import 'arquivos_io.dart'
    if (dart.library.js_interop) 'arquivos_web.dart';

// decisão em tempo de execução
final suportaBle = !kIsWeb &&
    (defaultTargetPlatform == TargetPlatform.android ||
     defaultTargetPlatform == TargetPlatform.iOS);
\end{codigo}
      \begin{conceito}[Por que importação condicional]
        \cod{dart:io} (arquivos, \emph{sockets}) não existe no navegador. Cada plataforma recebe o arquivo que sabe usar, com a mesma interface.
      \end{conceito}
    \end{column}
    \begin{column}{0.44\textwidth}
      \begin{definicao}[Política de mesma origem e CORS]
        O navegador bloqueia requisições de \cod{labmonitor.dev} para \cod{api.outro.dev}, a menos que a API responda com o cabeçalho \cod{Access-Control-Allow-Origin} autorizando a origem.
      \end{definicao}
      \begin{erro}[Funciona no celular, falha na web]
        Apps nativos não aplicam CORS; o navegador sim. A correção é \textbf{no servidor} (Módulo 17), não no app.
      \end{erro}
    \end{column}
  \end{columns}
\end{frame}

% =============================================================================
\section{Desktop e embarcados}

\begin{frame}[fragile]{Desktop e layout adaptativo}
  \begin{columns}[T]
    \begin{column}{0.56\textwidth}
\begin{terminal}
$ flutter create --platforms=windows,macos,linux .
$ flutter run -d linux
$ flutter build linux --release
\end{terminal}
\begin{codigo}[title=Um painel para cada largura]{Dart}
@override
Widget build(BuildContext context) => LayoutBuilder(
  builder: (context, c) => switch (c.maxWidth) {
    < 600 => const PainelCompacto(),    // lista
    < 840 => const PainelMedio(),       // grade
    _ => const PainelExpandido(),       // lista + detalhe
  },
);
\end{codigo}
    \end{column}
    \begin{column}{0.42\textwidth}
      \begin{conceito}[O que muda no desktop]
        \begin{itemize}
          \item Mouse: \emph{hover}, clique direito, rolagem.
          \item Teclado: atalhos (\cod{Shortcuts}, \cod{Actions}) e foco com Tab.
          \item Janelas redimensionáveis e menus.
          \item Instaladores próprios: MSIX, DMG, Snap ou DEB.
        \end{itemize}
      \end{conceito}
      \begin{dica}
        \cod{switch} com padrões relacionais (Dart 3) deixa as faixas de largura legíveis.
      \end{dica}
    \end{column}
  \end{columns}
\end{frame}

\begin{frame}[fragile]{Painel do laboratório num Raspberry Pi}
  \begin{columns}[T]
    \begin{column}{0.5\textwidth}
\begin{terminal}[]
# no computador de desenvolvimento
$ flutter pub global activate flutterpi_tool
$ flutterpi_tool build --arch=arm64 --release
# copiar build/flutter-pi para o Pi e rodar
$ flutter-pi --release ~/labmonitor
\end{terminal}
      \begin{conceito}[\cod{flutter-pi}]
        \emph{Embedder} leve que roda o Flutter direto na placa de vídeo do Raspberry Pi, sem ambiente gráfico: inicia rápido e ocupa a tela inteira, ideal para quiosques.
      \end{conceito}
    \end{column}
    \begin{column}{0.48\textwidth}
      \begin{exemplo}[LabMonitor em modo quiosque]
        Um Raspberry Pi 4 ou 5 com tela de toque fixado na parede mostra as bancadas, as reservas do dia e os alertas, recebendo dados por MQTT (Módulo 13).
      \end{exemplo}
      \begin{dica}[Hardware da placa]
        GPIO, I\textsuperscript{2}C e SPI podem ser acessados via FFI ou pacotes como \cod{dart\_periphery}: o painel lê sensores ligados diretamente ao Pi.
      \end{dica}
      \begin{historia}
        Montadoras já usam Flutter em centrais multimídia de carros, um dos usos mais conhecidos de embarcados.
      \end{historia}
    \end{column}
  \end{columns}
\end{frame}

% =============================================================================
\section{Fechamento}

\begin{frame}{Resumo}
  \begin{resumo}
    \begin{itemize}
      \item As abordagens vão do web/híbrido ao nativo; o Flutter compartilha tudo e desenha com motor próprio.
      \item Framework, engine e \emph{embedder} separam o Flutter das plataformas.
      \item Canais de plataforma trocam mensagens; Pigeon gera o código tipado; FFI chama C diretamente.
      \item Plugins federados e \emph{add-to-app} integram Flutter e código nativo nos dois sentidos.
      \item Web, desktop e embarcados reaproveitam o código, com layout adaptativo e atenção às diferenças (CORS, \cod{dart:io}).
    \end{itemize}
  \end{resumo}
\end{frame}

\begin{frame}{Glossário (1/2)}
  \begin{glossario}
    Multiplataforma & código escrito uma vez e executado em mais de um sistema \\
    JIT / AOT & compilação durante a execução / antes da distribuição \\
    \emph{Engine} & núcleo em C++ do Flutter: renderização, Dart \emph{runtime}, texto e canais \\
    \emph{Embedder} & camada que integra o \emph{engine} a uma plataforma específica \\
    Canal de plataforma & meio de troca de mensagens entre Dart e código nativo \\
    \emph{Codec} & regra que converte valores em bytes e de volta para atravessar o canal \\
    Pigeon & gerador de código tipado para os dois lados de um canal \\
    FFI & interface para chamar funções C a partir do Dart, compartilhando memória \\
  \end{glossario}
\end{frame}

\begin{frame}{Glossário (2/2)}
  \begin{glossario}
    Plugin federado & plugin dividido em API, interface de plataforma e implementações \\
    \emph{Add-to-app} & incorporação de telas Flutter a um app nativo existente \\
    PWA & site instalável, com manifesto e \emph{service worker} \\
    WebAssembly & formato binário executado pelos navegadores com desempenho próximo ao nativo \\
    CORS & mecanismo do navegador que controla requisições entre origens diferentes \\
    Layout adaptativo & layouts distintos por classe de tamanho de janela e forma de entrada \\
    Quiosque & aparelho dedicado a um único app, em tela cheia \\
    Matriz ponderada & técnica de decisão que soma notas multiplicadas por pesos \\
  \end{glossario}
\end{frame}

\begin{frame}{Questões de revisão}
  \begin{questoes}
    \item Compare React Native e Flutter quanto a quem desenha os componentes da tela. Que consequência isso tem para a aparência do app?
    \item Por que o Flutter compila AOT no iOS?
    \item Qual a diferença entre um canal de plataforma e o FFI? Dê um caso de uso para cada um.
    \item Que problema o Pigeon resolve em relação ao \cod{MethodChannel} escrito à mão?
    \item O que é um plugin federado e qual a vantagem dessa organização?
    \item Por que uma chamada HTTP pode funcionar no app Android e falhar no Flutter web?
    \item Quando o Flutter web \textbf{não} é uma boa escolha?
    \item Monte uma matriz ponderada para decidir entre Flutter e PWA para um painel de consulta do laboratório.
  \end{questoes}
\end{frame}

\begin{frame}{Exercícios}
  \begin{exercicio}[Exercício 1 --- Canal de plataforma]
    Crie um \cod{MethodChannel} que devolve o nome do modelo do aparelho no Android e no iOS, tratando a ausência de implementação.
  \end{exercicio}
  \begin{exercicio}[Exercício 2 --- FFI]
    Compile a função \cod{media\_movel} em C, chame-a pelo Dart e compare o tempo com uma versão em Dart para 10\,000 amostras.
  \end{exercicio}
  \begin{exercicio}[Exercício 3 --- Web e desktop]
    Gere as versões web e desktop do LabMonitor, tornando o painel adaptativo para as três classes de largura.
  \end{exercicio}
  \begin{projeto}[LabMonitor --- Etapa 19]
    Painel em modo quiosque (web ou Raspberry Pi) recebendo leituras por MQTT, e um canal de plataforma próprio para uma função nativa que nenhum plugin oferece.
  \end{projeto}
\end{frame}

\begin{frame}{Para aprofundar}
  \small
  \begin{itemize}
    \item Arquitetura do Flutter: \url{https://docs.flutter.dev/resources/architectural-overview}
    \item Canais de plataforma: \url{https://docs.flutter.dev/platform-integration/platform-channels}
    \item FFI: \url{https://dart.dev/interop/c-interop}
    \item \emph{Add-to-app}: \url{https://docs.flutter.dev/add-to-app}
    \item Flutter web: \url{https://docs.flutter.dev/platform-integration/web}
    \item flutter-pi: \url{https://github.com/ardera/flutter-pi}
  \end{itemize}
\end{frame}

\begin{frame}{Respostas das questões de revisão}
  \begin{respostas}
    \item O React Native usa os componentes nativos de cada sistema, que se parecem com os de outros apps; o Flutter desenha os próprios, com aparência idêntica nas plataformas e independente da versão do sistema.
    \item A App Store não permite gerar e executar código de máquina novo em tempo de execução; com AOT, o código já chega compilado.
    \item O canal troca mensagens serializadas e assíncronas com Kotlin/Swift (abrir uma API do sistema); o FFI chama funções C diretamente, inclusive de forma síncrona (filtro de sinais com muitas chamadas).
    \item Elimina nomes de métodos em texto e conversões manuais: a interface declarada gera código tipado nos dois lados, e erros aparecem na compilação.
    \item É um plugin dividido em API pública, interface comum e uma implementação por plataforma; novas plataformas podem ser adicionadas por pacotes separados.
    \item O navegador aplica a política de mesma origem: sem o cabeçalho CORS na resposta da API, a requisição é bloqueada.
    \item Em sites de conteúdo que dependem de SEO e de carregamento muito rápido da primeira página.
    \item Exemplo de critérios: acesso a BLE (peso 1, pois o painel só consulta), instalação sem loja (3), custo de manutenção (2), desempenho de gráficos (2). A PWA tende a vencer se o painel apenas exibe dados.
  \end{respostas}
\end{frame}

\encerramento{Módulo 20 --- Tópicos avançados e projeto final: IA, segurança, monetização e o projeto integrador.}

\end{document}
